\chapter{Literature Review}%
\label{chp:literature_review}

\section{Evolution of AI in Sports Analytics}
Computer vision (CV) has fundamentally altered the landscape of sports analytics, transitioning
from manual annotation to automated, high-fidelity measurement. Tracing the historical arc of
this evolution reveals a trajectory that began in the 1990s and early 2000s with manual
notational analysis. During this era, practitioners relied on painstaking frame-by-frame video
review and manual timing systems to derive basic performance metrics. The subsequent era of
``hand-crafted'' features saw the introduction of algorithms such as the Histogram of Oriented
Gradients (HOG), Scale-Invariant Feature Transform (SIFT), and various optical flow techniques.
While these early systems represented a step toward automation, they relied on heuristic-based
observation and were notoriously fragile, frequently failing under the non-uniform lighting and
high-entropy background noise characteristic of aquatic environments.

The paradigm shift occurred following the ``Deep Learning revolution,'' catalysed in large part
by the success of AlexNet in the 2012 ImageNet challenge~\cite{krizhevsky2012imagenet}. This
milestone demonstrated that hierarchical feature extraction through Deep Convolutional Neural
Networks (CNNs) could outperform any human-engineered descriptor. Building on the foundations
laid by earlier architectures such as LeNet~\cite{lecun2002gradient}, researchers began applying
deep learning to sports-specific tasks. In the mid-2010s, models such as
ResNet~\cite{he2016deep} became standard for robust feature representation, allowing for more
stable tracking in dynamic environments.

However, the adoption of these technologies has been uneven across the sporting landscape. Field
sports such as football, tennis, and athletics have received a disproportionate share of research
attention due to the relative ease of data collection in well-lit, terrestrial environments. In
contrast, swimming has lagged behind, primarily due to the unique optical challenges of the pool.
As noted by~\cite{Zhou2025}, AI-driven systems are now the standard for extracting kinematic
data and estimating pose in real time, but their application in swimming remains a frontier of
computer vision. The current paradigm is moving toward Foundation Models, defined as large-scale
models trained on massive, diverse datasets that exhibit zero-shot
generalisation~\cite{dosovitskiy2020image}. This thesis utilises models such as SAM~2 and
ViTPose, which provide a level of geometric and spatial reasoning, as well as temporal
consistency, previously unavailable in sports-specific contexts. Tools such as
DeepLabCut~\cite{mathis2018deeplabcut} and OpenPose~\cite{cao2019openpose} pioneered the
transfer of these capabilities to biomechanics, but the transition to the aquatic domain requires
a fundamental reassessment of how models perceive motion through water.

\section{An Adversarial Domain for Computer Vision}
While field sports such as football or athletics provide relatively stable visual environments
with predictable occlusions, the aquatic setting serves as an adversarial domain that
stress-tests the limits of image processing. Unlike terrestrial pose estimation, where the
medium of air is the standard, the water-air interface introduces non-linear geometric
distortions that violate the fundamental pinhole camera model
assumptions~\cite{agrawal2012theory}.

To understand the severity of this challenge, one must consider Snell's Law, which dictates that
$n_1 \sin \theta_1 = n_2 \sin \theta_2$. Because the refractive index of water ($n \approx
1.33$) differs significantly from air ($n \approx 1.00$), the apparent position of a swimmer's
joint underwater is displaced from its true anatomical
coordinate~\cite{harvey1998calibration}. Underwater camera systems require careful calibration
to compensate for this distortion, often necessitating specialised refractive camera models or
dedicated underwater calibration rigs~\cite{shortis2015calibration}. Without such compensation,
standard pose estimation models produce systematic errors in joint angle calculations, rendering
them insufficient for high-performance biomechanical analysis.

Beyond refraction, the environment is plagued by dynamic noise. As~\cite{giulietti_swimmernet_2023}
highlights, measurements are strongly influenced by disruptive elements such as bubbles, splashes,
and surface glare. These elements introduce severe, transient occlusions and
``caustics'', the dancing light patterns on the pool floor caused by the focusing effect of
surface waves. Indoor pool lighting often exacerbates these effects, creating high-contrast
reflections that can be mistaken for anatomical features by naive detectors. Furthermore,
swimming involves a unique ``multi-phase'' appearance problem: an athlete alternates between
above-water and below-water states. During a butterfly stroke recovery or a freestyle tumble
turn, the swimmer's body undergoes sudden, drastic changes in visual texture and limb visibility
that have no direct terrestrial analogue.

These environmental stressors explain why terrestrial benchmarks, such as
COCO~\cite{lin2014microsoft} or the MPII Human Pose dataset~\cite{andriluka20142d}, are largely
inadequate for the aquatic domain. While a model trained on COCO may excel at identifying a
pedestrian on a sidewalk, it lacks exposure to the blurred, refracted, and partially occluded
silhouettes of a submerged athlete. Einfalt et al.~\cite{einfalt2018activity} demonstrate that
even fine-tuned CNNs evaluated on real swimming footage exhibit systematic error modes in the
aquatic environment, confirming that models pre-trained on terrestrial data generalise poorly to
this domain. This necessitates the use of temporal memory mechanisms, such as the memory bank
implemented in SAM~2~\cite{ravi2024sam}, which allows a system to maintain object features even
when they are obscured by a splash or submerged beyond clear visibility. By treating the aquatic
environment as a structured noise problem rather than random interference, temporal dependencies
can be exploited to maintain tracking where classical optical flow methods fail systematically.

\section{From Handcrafted Features to Foundation Models}
Historically, structural feature detection in swimming relied on the Hough Transform or
colour-keying to isolate lane ropes or markers. However, these methods break down under varying
pool illumination and water turbulence~\cite{holzmann_measuring_2012}. Colour-keying, for
instance, is highly susceptible to the ``blue-shift'' of underwater light, whereby longer
wavelengths are absorbed more rapidly, causing markers to lose their distinct chromatic
signatures at depth. Similarly, classical optical flow algorithms such as Lucas-Kanade fail in
the presence of ``aperture problems'' caused by the repetitive, textureless nature of pool tiles
and the erratic motion of bubbles.

The subsequent shift toward CNNs, including U-Net and earlier You Only Look Once
(YOLO) variants, offered better robustness to noise. In the context of pose estimation, this era
saw the emergence of two primary strategies: ``bottom-up'' and ``top-down'' approaches.
Bottom-up models such as OpenPose~\cite{cao2019openpose} detect all body parts in an image and
then group them into individuals, which is computationally efficient but often struggles with the
limb-assignment ambiguities found in a crowded swimming lane. Top-down approaches first detect
the person and then perform pose estimation within that bounding box. For
instance,~\cite{al-majnoni_moar_2025} demonstrated the utility of YOLOv8 for athlete
identification, providing a stable crop for subsequent analysis. Nevertheless, CNNs are often
limited by their local receptive fields. Because a convolution operation only `sees' a small
neighbourhood of pixels at a time, it can struggle with the long-range spatial dependencies
required to identify a swimmer's full limb extension when the torso is clear but the hand is
obscured by a splash.

These limitations led to the adoption of Vision Transformers (ViT). Unlike convolutions, the
self-attention mechanism in ViT~\cite{dosovitskiy2020image} allows every token in the frame to
interact with every other token. Models such as ViTPose~\cite{xu2022vitpose} exploit this global
context to process the entire frame simultaneously, allowing the model to leverage the context of
the swimmer's torso to accurately predict the position of a wrist or ankle, even if that
specific joint is visually distorted by surface ripples. ViTPose is particularly notable for its
``plain ViT'' backbone and simple decoder head, an architectural choice that prioritises feature
richness over complex structural engineering and tends to improve robustness on
out-of-distribution data such as underwater footage.

The most recent evolution is the emergence of Foundation Models for segmentation and feature
extraction, notably the Segment Anything Model (SAM)~\cite{kirillov2023segment} and
DINO~\cite{caron2021emerging}. SAM~2 extends segmentation to the temporal domain through a
dedicated memory bank, enabling the segmentation of arbitrary objects such as underwater
calibration markers with zero-shot accuracy. These models offer strong geometric reasoning,
though not without trade-offs: the computational cost of global self-attention is considerably
higher than that of local convolutions, which presents challenges for real-time deployment. For
high-fidelity biomechanical research, however, where accuracy takes precedence over latency,
the capacity of Transformers to maintain a coherent kinematic chain through a turbulent
water-air interface represents a meaningful advance over the CNN era.

\section{The Challenge of Metric Scaling and Calibration}
A fundamental limitation in many CV-based swimming systems is the conversion of distance from
pixel space to real-world, metric units. A velocity of $x$ pixels per second is meaningless
without a reference scale. In sports biomechanics, the ``gold standard'' for this conversion is
3D multi-camera calibration, often using Direct Linear Transformation (DLT). As the constraints
of multi-camera acquisition in swimming demonstrate~\cite{sanders2015reliability}, poolside data
collection is limited by installation space, the requirement for waterproof housing, and the
difficulty of camera synchronisation, making this approach impractical in routine coaching
environments.

Alternative paradigms include wearable Inertial Measurement Units (IMUs), which can provide
direct acceleration data. However, IMUs face significant challenges in swimming, including data
contamination from water resistance and the fact that they are often prohibited in official
competitions. This leaves CV-based planar homography as the most viable alternative. Homography
assumes a mapping between two planes, for example the plane of the swimmer's lane and the camera
sensor. While~\cite{Tran2024} utilised lane ropes for this purpose, discrete, underwater markers
provide a more stable reference.

A critical issue in homography-based scaling is error propagation. Because velocity is a derived
quantity ($v = \Delta d / \Delta t$), a small error in marker detection or a slight violation of
the ``planar world'' assumption (for example, the swimmer moving deeper into the pool than the
calibration plane) can lead to significant inaccuracies in stroke length and speed calculations.
This project extends the logic of marker-based calibration by using SAM~2 to track these markers
with high temporal consistency. By establishing a fixed pixel-to-metre ratio based on known
marker distances, such as the 2.5\,m proposed in this thesis, the system achieves metric
accuracy without the overhead of an expensive sensor array or the instability of manual scaling.

\section{Addressing the Data Scarcity Gap: Model-in-the-Loop Annotation}
A persistent gap identified in the literature is the scarcity of publicly available, annotated
underwater datasets. To quantify this gap, one need only compare the scale of terrestrial
datasets such as COCO, which contains over 200{,}000 labelled
images~\cite{lin2014microsoft}, with aquatic datasets such as
SwimmerNET~\cite{giulietti_swimmernet_2023}, which are orders of magnitude smaller and often
lack specific features such as underwater markers. This scarcity is driven by ethical
restrictions on filming athletes, the logistical difficulty of underwater data collection, and a
lack of standardised annotation protocols for the aquatic domain.

A solution emerging in recent literature is Model-in-the-Loop (MITL) annotation, which sits
within the broader landscape of semi-supervised learning. Traditional semi-supervised techniques
such as pseudo-labelling~\cite{lee2013pseudo} and consistency regularisation (for example,
FixMatch~\cite{sohn2020fixmatch}) attempt to leverage unlabelled data by having a model predict
its own labels. MITL refines this by using high-capacity Foundation Models to generate
pseudo-labels. For example, by using SAM~2's zero-shot segmentation to identify markers in a
small subset of frames and propagating those masks through the temporal sequence, a high-quality
training set can be generated for a more efficient, task-specific detector such as YOLOv8.

This approach is further supported by self-supervised models such as
DINOv2~\cite{oquab2023dinov2}, which learn robust visual features without any labels at all.
However, MITL carries risks, the primary concern being ``confirmation bias,'' whereby the
student model (for example, YOLO) learns to replicate the systematic errors of the teacher model
(for example, SAM~2). This project mitigates such risks through human-in-the-loop verification,
ensuring that the automated MITL pipeline produces markers that are anatomically and
geometrically sound. This workflow represents a departure from traditional manual annotation and
directly addresses the research need for domain-specific underwater features.

\section{Summary of Research Needs}
The reviewed literature confirms a critical tension between the high-performance requirements of
swimming biomechanics and the adversarial nature of the aquatic environment. The transition from
manual notation to deep learning has provided the tools for automation, yet the field remains
hampered by three persistent challenges: the failure of traditional trackers during
high-turbulence occlusions; metric inaccuracy caused by refraction and poor calibration; and the
``cold start'' problem resulting from a lack of labelled underwater data.

This research bridges these gaps by synthesising a multi-stage pipeline that leverages recent
advances in Foundation Models. By using SAM~2 for stable temporal segmentation, the challenge of
occlusion and splash noise is addressed directly. By implementing a marker-based homography, a
robust solution to the metric scaling problem is provided that accounts for the refractive
properties of the medium. Finally, by employing an automated MITL pipeline, the work demonstrates
a scalable method for generating domain-specific data without the prohibitive cost of solely manual
annotation. Ultimately, this work moves beyond qualitative video analysis toward a metric-accurate
system that brings the rigour of terrestrial sports analytics to the complex, fluid environment
of competitive swimming.